\section{Exact observations and finite memory}\label{h:introduction}

An exact linear equation carries arbitrarily fine information. Indeed,
\(d\) independent Gaussian equations determine an unknown vector in
\(\R^d\) almost surely if all the equations can be retained. A learner
that reads the equations once and remembers only a finite state faces
a different problem. Increasing the required accuracy can force it to
read more equations even when every equation is noiseless.

\begin{definition}[Finite-state Gaussian regression]\label{h:model}
Fix integers \(d\ge2\), \(M\ge0\), and \(T\ge0\). A data-independent
shared seed may choose the learner's rules; the descriptions of its
transitions and output below are conditional on that seed. A signal
\(s\in S^{d-1}\) is chosen before sampling. At sample index
\(t\in\{1,\ldots,T\}\), the observation is
\[
 x_t\sim N(0,I_d),\qquad y_t=\ip{x_t}{s},
\]
where the rows \(x_1,\ldots,x_T\) are independent. A learner reads the
pairs once, in order. After each number of samples its persistent state
belongs to a prescribed set of at most \(2^M\) elements. The initial
state and shared seed are jointly independent of the signal and all
rows, but may depend on each other. Conditional on the seed and state,
a transition may use the entire current pair and fresh
randomness, with unrestricted computation, but only the next finite
state is retained.

The learner may stop, including before the first sample, but must stop
by the prescribed deterministic finite horizon \(T\). At a stopping index its unit
output \(\widehat s\in S^{d-1}\) is generated from the terminal state,
the stopping index, and fresh randomness only. In particular, a
discarded observation cannot be read again by the output rule.
The transitions and stopping rules may depend on \(d\), the requested
accuracy \(\epsilon\), and the index, but not on \(s\).

The experiment must be jointly measurable in the signal, the complete
pre-generated rows, the shared seed, the initial state, and all fresh
transition and output randomness. Conditional on the seed, each fixed
state/index transition and stopping probability must be Borel in
\((x,y)\), or measurable in the completion of
\(\lambda_1(dx,dy)=\gamma_1(dx)\,dy\), where \(\gamma_1\) is standard
Gaussian probability on \(\R^d\). Fresh randomness supplies no
information about the signal or samples beyond the permitted rule
arguments. Accuracy means
\[
 \arccos\ip{\widehat s}{s}\le\epsilon.
\]
\end{definition}

The state bound measures bits, not real registers. The rules can be
nonuniform in the dimension and sample index, and randomization is
allowed both in choosing the rules and during the run. The model fixes
the rows in advance through their independent Gaussian law: it does
not allow a learner to choose a row or revisit an earlier one.

Write \(\sigma_d\) for uniform probability on \(S^{d-1}\). Our main
statement assumes success only after averaging over this prior.

\begin{theorem}[Memory and precision]\label{h:main}
For every fixed finite \(A>0\), there are constants \(c_A>0\) and
\(d_A\) such that the following holds. Let \(d\ge d_A\), let
\(0\le M\le Ad^2\) be an integer, and let \(0<\epsilon\le1/10\).
For any learner in Definition~\ref{h:model}, draw
\(S\sim\sigma_d\) independently of the rows and the learner's
data-independent randomness. If
\[
 \Prob\{\arccos\ip{\widehat s}{S}\le\epsilon\}\ge\frac23,
\]
then its prescribed horizon satisfies
\[
 T\ge c_A\,d\log(1/\epsilon).
\]
\end{theorem}

Here and throughout the paper logarithms are natural unless a base is
shown. The constants in Theorem~\ref{h:main} may depend on the fixed
number \(A\). The statement makes no uniform claim when \(A=A(d)\)
tends to infinity. No relation between \(\epsilon\) and \(M\) is
assumed.

\begin{corollary}[Subquadratic memory]\label{h:subquadratic}
Let \(M(d)\ge0\) be integer-valued with \(M(d)=o(d^2)\), and let
\(0<\epsilon(d)\le1/10\) be any accuracy sequence. A family of learners
in Definition~\ref{h:model} with uniform-sphere average success at
least \(2/3\) satisfies
\[
 T(d)\ge c\,d\log(1/\epsilon(d))
\]
for all sufficiently large \(d\), where \(c>0\) is absolute. The
dimension threshold may depend on the memory sequence.
\end{corollary}

\begin{corollary}[A guarantee for every signal]\label{h:worst-signal}
The conclusions of Theorem~\ref{h:main} and
Corollary~\ref{h:subquadratic} also hold when their average-success
assumption is replaced by
\[
 \Prob\{\arccos\ip{\widehat s}{s}\le\epsilon\}\ge\frac23
 \qquad\text{for every }s\in S^{d-1},
\]
where this probability is over the rows and the learner's randomness.
\end{corollary}

The every-signal consequence follows by averaging the same learner
over \(\sigma_d\); it does not say that every individual signal is hard
for every learner. We prove Theorem~\ref{h:main} and both corollaries
after completing the backward argument in Section~\ref{h:end:section}.

\paragraph{The original experiment and measurable versions.}
All success probabilities refer to the original jointly measurable
experiment. Joint independence of the initial state and seed from the
signal and rows permits conditioning on the seed without changing the
prior or supplying information through initialization.
Section~\ref{h:stream:section} replaces completed-measurable rules
by Borel versions separately for almost every seed, preserving the
corresponding prior success probabilities. Their average is taken
in the original experiment; no jointly measurable choice of all these
versions is assumed.

\subsection{The precision problem and its predecessors}

Memory constraints in statistical learning were studied before this
precision question. Shamir's finite-message framework included
bounded-memory online processing as a special case
\cite[Section~2, Definition~1]{Shamir2014}. In its general form a
protocol may retain an entire transcript, so its per-message budget is
not a bound on total persistent memory. Steinhardt and Duchi proved
memory-dependent minimax risk bounds for sparse noisy regression
\cite[Section~1.1, Equation~(2)]{SteinhardtDuchi2015}. Their hard
family uses positive noise and a different covariate distribution.
These works established that information retained from earlier samples
can be a statistical resource in its own right.

Steinhardt, Valiant, and Wager connected memory, communication, and
statistical queries, and formulated a quadratic-memory versus
exponential-sample question for parity learning
\cite[Section~1.1, Conjecture~1]{SVW2016}. Raz proved the qualitative
separation for parity \cite[Theorem~1 and Remark~1.1]{Raz2016} and
developed a spectral method for a wider class of learning problems
\cite{Raz2017}. The qualifier \emph{qualitative} matters: Raz's
theorem and the published conjecture have different numerical
memory coefficients. These discrete results explain the usefulness of
finite-width computation models, but they do not supply an estimate
for exact real-valued observations.

Among the sources discussed here, the closest continuous-regression
comparison is due to Sharan, Sidford, and Valiant
\cite{SSV2019}. Their Theorem~1 uses independent \(N(0,I_d)\)
covariates, a fixed unit target, and independent uniform additive noise
of half-width \(2^{-d/5}\). With at most \(d^2/4\) bits and Euclidean
accuracy \(\epsilon=d^{-r}\), it gives an
\(\Omega(d\log r)\) sample lower bound for \(r\) from a sufficiently
large absolute constant up to order \(d/\log d\), in the range stated
there. Their proof experiment draws the target uniformly on the sphere
and averages over it and the samples \cite[Section~2]{SSV2019}.
Their nonuniform finite-width layered program is a predecessor of the
model above; the measurable randomized rules and bounded stopping used
here are handled directly in Section~\ref{h:stream:section}.

The observation law is the important difference. A lower bound with
additive noise does not by itself bound the more informative exact
experiment. Theorem~\ref{h:main} treats the exact-observation precision
question for isotropic Gaussian rows. Section~1.1 of \cite{SSV2019}
asks whether the first-order sample dependence on precision is optimal
under bounded memory. This is distinct from its separate
condition-number conjecture for ill-conditioned Gaussian laws and
constant Euclidean accuracy in \cite[Section~1.1]{SSV2019}.

Dagan, Kur, and Shamir obtained related quadratic-space obstructions
for two other linear-prediction tasks \cite{DKS2019}. Their published
Theorem~3 asks for a unit approximate null vector of independent
Gaussian rows. Their Theorem~10 asks for small empirical residual on
a finite random-order system built from jointly conditioned sphere
rows and one fixed equation. These are different targets and
observation laws from the fixed-signal stream considered here.

\paragraph{A consequence for the noisy experiment.}
Theorem~\ref{h:main} implies an
\(\Omega_A(d r\log d)\) sample lower bound in the noisy Gaussian
experiment of Sharan, Sidford, and Valiant \cite{SSV2019}, for every
fixed finite \(A>0\), memory \(M\le Ad^2\), and Euclidean accuracy
\(e=d^{-r}\) in their fine-accuracy range described above,
for sufficiently large \(d\). Here success is at least \(2/3\)
averaged over the uniform-sphere target, and learners retain the state,
deterministic-horizon, randomness, and joint and per-seed measurability
conventions of Definition~\ref{h:model}, with a terminal estimate
\(V\in\R^d\) depending only on the terminal state, stopping index,
and fresh randomness. An exact-observation learner generates the
prescribed independent uniform noise inside each current-pair
transition and simulates the noisy learner with the same states and
horizon. For completed rules, the noisy pair law has a density with
respect to \(\lambda_1\), so the Borel-version and seed-conditioning
argument of Section~\ref{h:stream:section} applies before this
simulation. On \(\|V-S\|\le e\le1/20\), normalizing \(V\) gives
angular error at most \(\arcsin e\le2e\); assign any unit output when
\(V=0\). Theorem~\ref{h:main} at accuracy \(2e\) therefore gives
\(T\ge c_A d\log(1/(2e))=\Omega_A(d r\log d)\).

\subsection{Why the precision scale is natural}

With real registers, the classical Kaczmarz projection step
\cite{Kaczmarz1937}, in the form stated by Strohmer and Vershynin
\cite[Algorithm~1]{SV2009}, gives a simple benchmark. Starting with
\(z_0=0\), use
\begin{equation}\label{h:kaczmarz}
 z_t=z_{t-1}+
 \frac{y_t-\ip{x_t}{z_{t-1}}}{\norm{x_t}^2}\,x_t
 \qquad(x_t\ne0),
\end{equation}
and set \(z_t=z_{t-1}\) on the Gaussian-null event \(x_t=0\).
For \(u_t=x_t/\|x_t\|\), the error is multiplied by
\(I-u_tu_t^{\mathsf T}\). The direction \(u_t\) is fresh and uniform
on the sphere, so \(\E(u_tu_t^{\mathsf T})=I_d/d\). Conditional on
the previous error,
\[
 \E\!\left(\norm{z_t-s}^2\mid z_{t-1}\right)
  =(1-1/d)\norm{z_{t-1}-s}^2.
\]
Induction gives
\begin{equation}\label{h:kaczmarz-contraction}
 \E\norm{z_t-s}^2=(1-1/d)^t.
\end{equation}
This calculation explains the scale \(d\log(1/\epsilon)\) for fresh
isotropic observations. Strohmer and Vershynin's theorem concerns
norm-weighted sampling from a fixed finite matrix; the equality above
is the separate symmetry calculation for this stream. The displayed
update stores real scalars and supplies no finite-bit upper bound.

\subsection{A backward proof under the original prior}

Conditioning the uniform prior on one raw exact observation almost
surely confines the signal to a hyperplane section of the sphere. That
conditional law is singular relative to the original spherical
measure. We instead keep the original prior fixed and study the
probability of future success from a prescribed state.

Start the learner at a chosen state and index, whether or not the state
was reached in an actual run. Let \(h(s)\) be its success probability
with signal \(s\), using fresh future samples. For a closed Euclidean
ball \(B(z,r)\), our invariant is
\begin{equation}\label{h:intro:invariant}
 \int_{B(z,r)}h\,d\sigma_d\le(Rr)^{(d-1)/2}
 \qquad(z\in S^{d-1},\ r>0).
\end{equation}
The number \(R\) bounds the concentration of successful signals; it
is not a radius supporting a posterior. A terminal output rule has
\(R=\epsilon\). Moving backward across a block of order \(d\)
observations multiplies \(R\) by at most a constant for fixed
\(M/d^2\). Constant average success requires \(R\) bounded below, so
reaching it from \(\epsilon\) takes order \(\log(1/\epsilon)\) blocks.

The bound at every radius makes this iteration close. Restrict
\(h\,d\sigma_d\) to a ball \(B(z,r)\), without conditioning on reaching
a state or dividing by the restricted mass. This finite measure has
an ambient ball bound \(H t^\beta\), with \(\beta=(d-1)/2\) and
\(H\) controlled by \(R^\beta\). The general projection estimate in
Section~\ref{h:projection} gives the joint law of a Gaussian matrix
and its exact labels an \(L^q\) density whose norm is bounded by
\[
 e^{Cd}H r^{\beta-m(1-1/q)}
\]
for a block of \(m\le q\) rows, where \(q\) is proportional to \(d\).
The density loses a radius power. The volume of the common projected
ellipsoid is at most \(e^{Cd}r^m\), and its power \(1-1/q\) restores
exactly that loss in H\"older's inequality.

The geometric step controls inverse distances to affine spans, hence
an inverse Gram determinant. Local growth and inverse-distance
averaging belong to the classical projection methods represented by
Mattila \cite[Lemma~5.1 and Theorem~5.4]{Mattila1975}; the
independent-copy expansion and successive orthogonalization have a
closer regression predecessor in
\cite[Section~7, Lemma~12]{SSV2019}. We derive the exact multirow
density and its quantitative bounds here. Its finite-measure
formulation does not require a density relative to spherical area.

Section~\ref{h:blocks} combines the projection and volume bounds
with a second application of H\"older on the \(N\) possible
destinations. Selection costs \(N^{1/q}\), so taking the
\((d-1)/2\)-th root in \eqref{h:intro:invariant} gives the memory
contribution \(\exp(O(\log N/d^2))\).
Section~\ref{h:stream:section} constructs fresh continuation
functions, counts only stopping indices within one block, and passes
from Borel rules to the full randomized model. Reflection across the
span of the observed rows gives the independent constant-accuracy
endpoint in Section~\ref{h:end:section}, completing the theorem.

The general projection estimate also applies to a lower-dimensional
cube prior on the sphere. That application has its own distributional
precision bound; the conclusion throughout the stated accuracy range
uses an every-signal guarantee and the separate uniform-sphere
endpoint. The appendices give distinct exact-label interfaces,
compare coefficient normalizations through one backward bound,
and retain a sharper full-data residual estimate.
